\begin{frame}{例子：一个看起来完整但其实不可信的实验}
  \Core{实验方法课必须先教学生识别坏实验，否则工具越熟练，错误可能越系统化。}
  \begin{block}{表面上很完整}
    有吞吐图、有延迟表、甚至有消融，但 workload 来源和对照配置没有交代清楚。
  \end{block}
  \pause
  \begin{block}{真正的问题}
    结果无法回答“为什么变好”以及“在哪些条件下不再成立”。
  \end{block}
  \pause
  \begin{block}{课程结论}
    一个实验是否可信，不取决于图表数量，而取决于因果链和边界是否被写清。
  \end{block}
\end{frame}

\begin{frame}{例子：端到端数字与探针数字不一致时怎么办}
  \Core{最有研究价值的情况往往不是数字全都一致，而是局部证据与端到端现象出现张力。}
  \begin{block}{情况一}
    GPU 利用率更高，但 TTFT 更差，说明等待或同步代价可能在放大。
  \end{block}
  \pause
  \begin{block}{情况二}
    KV 命中率更高，但 goodput 没提升，说明命中本身未必落在主导瓶颈上。
  \end{block}
  \pause
  \begin{block}{教学意义}
    探针不是为了给端到端结果背书，而是为了帮助发现矛盾并修正假设。
  \end{block}
\end{frame}

\begin{frame}{代码例子：最小 benchmark 模板应包含什么}
  \Core{一份合格的 benchmark，不只是能跑，还应显式携带 workload 和环境信息。}
  \begin{block}{代码片段}
    {\ttfamily\small
    \detokenize{seed(0)}\par
    \detokenize{warmup_requests = 10}\par
    \detokenize{num_requests = 100}\par
    \detokenize{measure(ttft, throughput, p99)}\par
    }
  \end{block}
  \pause
  \begin{block}{课堂应追问}
    哪些配置是影响 workload 定义的，哪些配置只是运行时细节？
  \end{block}
  \pause
  \begin{block}{系统含义}
    benchmark 模板的显式化，本身就在减少“实验只存在于一台机器上”的偶然性。
  \end{block}
\end{frame}

\begin{frame}{代码例子：最小日志汇总脚本}
  \Core{实验课需要让学生看到，图表往往来自对原始日志的二次组织。}
  \begin{block}{代码片段}
    {\ttfamily\small
    \detokenize{records = load_jsonl("run.log") }\par
    \detokenize{p99 = percentile([r["latency_ms"] for r in records], 99)}\par
    \detokenize{goodput = sum(r["ok"] for r in records) / total_time}\par
    }
  \end{block}
  \pause
  \begin{block}{课堂应追问}
    如果日志字段设计不对，后面哪些指标将根本算不出来？
  \end{block}
  \pause
  \begin{block}{课程结论}
    可验证性不只发生在论文写作阶段，而是从日志字段设计就已经开始了。
  \end{block}
\end{frame}

\begin{frame}{例子：消融实验最容易做假的地方}
  \Core{消融的真正任务是隔离变量，而不是制造更多子图。}
  \begin{block}{伪消融一}
    一次改了多个控制点，却只关掉其中一个开关就说完成了消融。
  \end{block}
  \pause
  \begin{block}{伪消融二}
    workload 或参数同时变化，导致数字变化无法归因。
  \end{block}
  \pause
  \begin{block}{课程要求}
    每一项消融都要明确写出“被隔离的变量”和“预期影响的指标”。
  \end{block}
\end{frame}

\begin{frame}{课堂练习：写一页可信实验卡片}
  \Core{实验方法训练最好用一页纸完成：问题、对照、指标、失败条件。}
  \begin{block}{练习一}
    写出你的 workload 定义，至少包含输入长度分布、输出长度分布和到达模式。
  \end{block}
  \pause
  \begin{block}{练习二}
    写出一个主指标和两个辅助探针，并说明三者的因果关系假设。
  \end{block}
  \pause
  \begin{block}{练习三}
    用一句话写出：什么结果出现时，你会承认当前机制没有兑现预期收益？
  \end{block}
\end{frame}
